\documentclass[11pt,twoside]{book}
\usepackage[T1]{fontenc}
\usepackage[utf8]{inputenc}
\usepackage{libertine}
\usepackage[margin=1in]{geometry}
\usepackage{amsmath,amssymb,amsthm,mathtools}
\usepackage{booktabs}
\usepackage{enumitem}
\usepackage{xcolor}
\usepackage[hidelinks]{hyperref}
\usepackage[nameinlink,capitalise]{cleveref}
\definecolor{scigreen}{HTML}{356859}
\newtheorem{theorem}{Theorem}[chapter]
\newtheorem{definition}[theorem]{Definition}
\newcommand{\vect}[1]{\boldsymbol{#1}}
\begin{document}
\begin{titlepage}
\centering
\vspace*{0.85in}
{\Huge\bfseries Notes on Structured Inquiry\par}
\vspace{15pt}{\Large A scientific-style book starter\par}
\vspace{20pt}\textcolor{scigreen}{\rule{0.72\textwidth}{2pt}}\par
\vfill
{\Large Author Name\par}\vspace{6pt}{\large Research group or institution\par}
\vspace{0.5in}{\small Version 1.0 \textbullet{} \the\year\par}
\end{titlepage}
\frontmatter
\chapter*{Overview}
This starter foregrounds a navigable table of contents, labelled results,
cross-references, compact tables, and appendices. Replace the examples with
the argument and evidence of your own scientific monograph.
\tableofcontents
\mainmatter
\chapter{Questions and definitions}\label{ch:questions}
\section{A precise vocabulary}
\begin{definition}
An observation is a recorded statement about a system together with the
conditions under which it was recorded.
\end{definition}
Good definitions reduce ambiguity before a calculation begins. The workflow
in \cref{ch:methods} uses this distinction to separate observations from
interpretations.
\chapter{Methods and evidence}\label{ch:methods}
\section{A small reproducible protocol}
\begin{enumerate}[label=\arabic*.]
  \item State the question and the measurable outcome.
  \item Record inputs, instruments, and environmental conditions.
  \item Preserve the raw observations before summarising them.
\end{enumerate}
\begin{theorem}[A transparent estimate]\label{thm:estimate}
If each observation is recorded with its unit and uncertainty, a later reader
can reproduce the reported estimate from the same evidence.
\end{theorem}
\begin{proof}
The unit identifies the scale of the quantity, while the uncertainty records
the precision claimed by the measurement. Together they determine the
interpretation of the estimate.
\end{proof}
\section{Example results}
\begin{table}[ht]
\centering
\caption{Illustrative measurements}\label{tab:measurements}
\begin{tabular}{lrrr}
\toprule Sample & Mean & Uncertainty & Unit\\\midrule
A & 12.4 & 0.3 & mm\\ B & 11.8 & 0.2 & mm\\ C & 12.1 & 0.4 & mm\\
\bottomrule
\end{tabular}
\end{table}
Table~\ref{tab:measurements} demonstrates a compact result layout.
\chapter{Interpretation}
\section{From result to claim}
A result is not automatically an explanation. Make the inferential step
explicit, identify competing explanations, and state what future observation
would distinguish them.
\backmatter
\appendix
\chapter{Notation and conventions}
Use appendices for derivations, supplementary tables, code listings, or a
notation index that would interrupt the main argument.
\chapter*{Acknowledgements}
Replace this neutral page with acknowledgements and contribution statements.
\end{document}
